\documentclass[11pt,a4paper]{article}

\usepackage{amsmath,amssymb,amsthm}
\usepackage{geometry}
\usepackage{hyperref}
\usepackage{booktabs}
\usepackage{listings}
\usepackage{xcolor}
\usepackage{parskip}
\usepackage{microtype}

\geometry{margin=2.5cm}

\hypersetup{
  colorlinks=true,
  linkcolor=blue!60!black,
  urlcolor=blue!60!black,
  citecolor=blue!60!black
}

\lstset{
  basicstyle=\ttfamily\small,
  backgroundcolor=\color{gray!10},
  frame=single,
  framerule=0pt,
  breaklines=true
}

\theoremstyle{plain}
\newtheorem{theorem}{Theorem}[section]
\theoremstyle{remark}
\newtheorem{remark}{Remark}[section]

% Convenience macros
\newcommand{\vth}{\vartheta}
\newcommand{\vths}{\vartheta^*}
\newcommand{\zth}{z_{\vth}}
\newcommand{\zom}{z_{\omega}}
\newcommand{\zdel}{z_{\delta}}
\newcommand{\clip}{\mathrm{clip}}
\newcommand{\dmax}{\delta_{\max}}
\newcommand{\dcmd}{\delta_{\mathrm{cmd}}}
\newcommand{\Fmax}{F_{\max}}

\title{\textbf{Full-System Backstepping Landing Control}\\
\large Mathematical Derivation and Stability Analysis}
\author{}
\date{}

\begin{document}
\maketitle
\tableofcontents
\newpage

% =============================================================================
\section*{Abstract}
% =============================================================================

This document derives and formally analyses a \textbf{full-system backstepping} landing controller for a planar thrust-vector-controlled (TVC) rocket. A single composite Lyapunov function simultaneously covers position, attitude, and the nozzle actuator. The derivation proceeds in four backstepping steps, producing virtual controls for desired pitch angle, angular rate, and nozzle deflection, before arriving at the real nozzle command. Global practical stability under actuator saturation is proved by establishing that the closed-loop system is \textbf{input-to-state stable (ISS)} with respect to actuator clipping errors. ISS is a property we prove for our specific system --- not an assumed result --- by constructing a Lyapunov inequality of the form $\dot V \leq -c\|z\|^2 + D'$, where $D'$ depends on the magnitude of the clipping errors. This guarantees the state remains bounded whenever the disturbances are bounded. The full proof with explicit quantitative bounds is in Section~\ref{sec:iss}.

% =============================================================================
\section{Plant Description}
% =============================================================================

The system is a planar TVC rocket in the vertical plane. A single engine deflects by angle $\delta$ from the body axis, and $\delta$ is governed by a first-order actuator. The rocket operates in the low-speed aerodynamic regime, so aerodynamic drag, normal force, and a pitching moment act on it. All aerodynamic coefficients are assumed known.

\subsection*{State vector and control inputs}

\[
  s = [x,\ y,\ \dot x,\ \dot y,\ \vth,\ \dot\vth,\ \delta]^T \in \mathbb{R}^7
\]

Control inputs: $\sigma \in [\sigma_{\min}, 1]$ (throttle) and $\dcmd \in [-\delta_{\max,\mathrm{cmd}},\, \delta_{\max,\mathrm{cmd}}]$ (nozzle command).

\subsection*{System parameters}

\begin{center}
\begin{tabular}{lll}
\toprule
Symbol & Meaning & Units \\
\midrule
$m$           & Total mass (constant)                         & kg \\
$J$           & Moment of inertia about CoM (constant)        & kg\,m$^2$ \\
$l_{cp}$      & Distance from CoM to nozzle exit             & m \\
$g$           & Gravitational acceleration, 9.81              & m/s$^2$ \\
$\Fmax$       & Maximum thrust                                & N \\
$\dmax$       & Physical nozzle deflection limit              & rad \\
$\tau_\delta$ & Nozzle actuator time constant                 & s \\
$\rho$        & Air density, 1.225                            & kg/m$^3$ \\
$S_m$         & Reference cross-sectional area                & m$^2$ \\
$l$           & Reference length                              & m \\
$C_x$         & Axial drag coefficient, 0.358                 & --- \\
$C_{y\alpha}$ & Normal force coefficient, 0.05403             & deg$^{-1}$ \\
$C_{m\alpha}$ & Pitching moment coefficient, 1.054            & rad$^{-1}$ \\
\bottomrule
\end{tabular}
\end{center}

% =============================================================================
\section{Simplifying Assumptions}
% =============================================================================

\textbf{A1. Constant mass.} $\dot m = 0$, $\dot J = 0$, $\dot l_{cp} = 0$.

\textit{In a real rocket, fuel is consumed during the burn, reducing mass and shifting the centre of mass. Here we assume the landing burn is short enough that mass change is negligible --- a standard assumption for terminal descent from a few hundred metres altitude.}

\textbf{A2. Low-speed aerodynamic regime.} $V \leq 100$ m/s (Mach $< 0.3$): $C_x = 0.358$, $C_{m\alpha} = 1.054$ rad$^{-1}$, $C_{y\alpha} = 0.05403$ deg$^{-1}$.

\textit{Below Mach 0.3 compressibility effects are small and aerodynamic coefficients are approximately constant. This allows us to use fixed coefficient values throughout the landing manoeuvre rather than scheduling them with Mach number.}

\textbf{A3. All parameters known.} $C_{m\alpha}$, $C_x$, $C_{y\alpha}$, $J$, $l_{cp}$, $\tau_\delta$ are all known.

\textbf{A4. Small nozzle deflection for translational linearisation.} $|\delta| \leq \dmax \leq 0.25$ rad, so:
\begin{align}
  \sin(\vth + \delta) &\approx \sin\vth + \delta\cos\vth \\
  \cos(\vth + \delta) &\approx \cos\vth - \delta\sin\vth
\end{align}
Error: $|\sin(\vth+\delta) - (\sin\vth + \delta\cos\vth)| \leq \delta^2/2 < 3\%$ for $|\delta| \leq 0.25$ rad.

\textit{The nozzle angle is physically limited to a narrow range (here $\pm 0.25$ rad $\approx \pm 14^{\circ}$) by the gimbal mechanism. Because $\delta$ is small, the first-order Taylor expansion of $\sin(\vth+\delta)$ and $\cos(\vth+\delta)$ around $\delta=0$ is accurate to within 3\%. This separates the dominant thrust direction (governed by body pitch $\vth$) from the small torque contribution of the nozzle deflection $\delta$, making the equations tractable for backstepping.}

\textbf{A5. Small pitch error for rotational linearisation.} $|\zth| = |\vth - \vths| \ll 1$ rad during the controlled transient, so $\sin(\vths + \zth) \approx \sin\vths + \cos\vths\cdot \zth$. This allows the position dynamics to be written in terms of $\zth$ explicitly, enabling a single Lyapunov function.

\textit{The backstepping design requires writing the time derivative of $V_\mathrm{pos}$ as an explicit function of the pitch error $\zth$. This is only possible if we can linearise $\sin\vth$ around the desired pitch $\vths$. The assumption holds whenever the attitude controller is fast enough to keep the tracking error small --- which is enforced by the inner-loop gains $k_\vth$ and $k_\omega$. It is the only assumption that creates a conservatism gap: if $|\zth|$ ever grows large (e.g., at the very start with a $30^{\circ}$ initial tilt), the linearisation error acts as an additional bounded disturbance absorbed into the ISS bound.}

\textbf{A6. Minimum throttle.} $\sigma \geq \sigma_{\min} > 0$, so $g_2 = F l_{cp}/J > 0$ always.

\textit{The control authority of the nozzle enters the rotational dynamics as $g_2 = F l_{cp}/J$. If the engine were shut off ($\sigma = 0$), $g_2 = 0$ and the virtual control $\alpha_2$ in Step 3 would involve division by zero (see equation (VC2)). Enforcing a minimum throttle $\sigma_{\min}$ keeps $g_2$ bounded away from zero and the controller well-defined at all times.}

\textbf{A7. Full state measurement.} $(x, y, \dot x, \dot y, \vth, \dot\vth, \delta)$ is measurable without noise.

% =============================================================================
\section{Equations of Motion}
% =============================================================================

\subsection{Full Nonlinear Equations}

The thrust vector makes angle $(\vth + \delta)$ from the vertical in the inertial frame. Projecting thrust and body-frame aerodynamic forces ($X_b$ axial drag, $Y_b$ normal force) onto the inertial axes, and including the nozzle actuator and pitch moment, the exact equations of motion are:

\begin{align}
  m\ddot x &= F\sin(\vth + \delta) + X_b\sin\vth + Y_b\cos\vth \tag{NL1} \\
  m\ddot y &= F\cos(\vth + \delta) - mg + X_b\cos\vth - Y_b\sin\vth \tag{NL2} \\
  J\ddot\vth &= F\,l_{cp}\,\sin\delta + C_{m\alpha}\,\alpha\,q_\infty S_m\,l \tag{NL3} \\
  \tau_\delta\,\dot\delta &= -\delta + \dcmd \tag{NL4}
\end{align}

with $F = \sigma \Fmax$, $X_b = -C_x q_\infty S_m$, $Y_b = C_{y\alpha}\,\alpha\,q_\infty S_m$, $q_\infty = \tfrac{1}{2}\rho V^2$, $V = \sqrt{\dot x^2 + \dot y^2}$, and angle of attack:
\[
  \alpha = \vth - \mathrm{atan2}(\dot x,\, \dot y)
\]

Note that (NL1)--(NL2) are fully coupled: the translational accelerations depend on the nozzle angle $\delta$ through $\sin(\vth+\delta)$ and $\cos(\vth+\delta)$, and on pitch $\vth$ through the aerodynamic projections. The rotational equation (NL3) is driven by $\delta$ through the moment arm $l_{cp}$.

\subsection{Small-\texorpdfstring{$\delta$}{delta} Linearisation (Assumption A4)}

Applying A4 to (NL1)--(NL2):

\begin{align}
  m\ddot x &= F\sin\vth + F\delta\cos\vth + X_b\sin\vth + Y_b\cos\vth \tag{T1} \\
  m\ddot y &= F\cos\vth - mg - F\delta\sin\vth + X_b\cos\vth - Y_b\sin\vth \tag{T2}
\end{align}

with $F = \sigma\Fmax$, $X_b = -C_x q_\infty S_m$, $Y_b = C_{y\alpha}\,\alpha\,q_\infty S_m$, $q_\infty = \tfrac{1}{2}\rho V^2$, and angle of attack $\alpha = \vth - \mathrm{atan2}(\dot x, \dot y)$.

This is the angle between the body axis and the velocity vector, measured positive nose-up. At zero airspeed ($V < V_{\min} \approx 0.1$ m/s), aerodynamic forces vanish and $\alpha$ is set to $\vth$ by convention.

\textbf{Structural observation.} In (T1)--(T2), $\vth$ drives position through the dominant terms $F\sin\vth$, $F\cos\vth$, while $\delta$ influences position through the smaller coupling terms $F\delta\cos\vth$, $-F\delta\sin\vth$. Retaining these coupling terms is what makes the full-backstepping approach stronger than a simple outer-inner loop design.

Applying A4 to the rotational equation ($\sin\delta \approx \delta$) and keeping the actuator equation exact:
\begin{align}
  J\ddot\vth &= F\,l_{cp}\,\delta + C_{m\alpha}\,\alpha\,q_\infty S_m\,l \tag{R1} \\
  \tau_\delta\,\dot\delta &= -\delta + \dcmd \tag{A1}
\end{align}

Define the shorthand:
\begin{align}
  f_2 &= \frac{C_{m\alpha}\,\alpha\,q_\infty S_m\,l}{J} \\
  g_2 &= \frac{F\,l_{cp}}{J} > 0
\end{align}

so that $\ddot\vth = g_2\,\delta + f_2$. Note that $C_{m\alpha} > 0$ means the rocket body is \textbf{aerodynamically unstable} --- a positive angle of attack produces a nose-up pitching moment, which amplifies the perturbation. The controller must actively cancel $f_2$ via the (VC2) term $+\dot\alpha_1 - f_2$ in the nozzle reference.

% =============================================================================
\section{Control Objective and Error Coordinates}
% =============================================================================

Landing target:
\[
  s^* = (x_d,\ 0,\ 0,\ 0,\ 0,\ 0,\ 0)^T
\]

Error coordinates:
\begin{align}
  e_x &= x - x_d, \quad e_y = y \\
  \dot e_x &= \dot x, \quad \dot e_y = \dot y
\end{align}

The full objective: $(e_x,\, e_y,\, \dot x,\, \dot y,\, \vth,\, \dot\vth,\, \delta) \to 0$ as $t \to \infty$.

% =============================================================================
\section{Unified Lyapunov Function --- Structure and Motivation}
% =============================================================================

The full-backstepping Lyapunov function is built up one term at a time --- one $\tfrac{1}{2}z_i^2$ per step --- until the entire state is covered:

\begin{equation}
  V = \underbrace{\tfrac{1}{2}k_{px}e_x^2 + \tfrac{1}{2}\dot x^2 + \tfrac{1}{2}k_{py}e_y^2 + \tfrac{1}{2}\dot y^2}_{V_{\mathrm{pos}}} + \underbrace{\tfrac{1}{2}\zth^2}_{\text{Step 2}} + \underbrace{\tfrac{1}{2}\zom^2}_{\text{Step 3}} + \underbrace{\tfrac{1}{2}\zdel^2}_{\text{Step 4}}
  \tag{V}
\end{equation}

where the backstepping errors are defined sequentially:
\[
  \zth = \vth - \vths, \quad \zom = \dot\vth - \alpha_1, \quad \zdel = \delta - \alpha_2
\]

\begin{itemize}
  \item $\zth$: how far the actual pitch $\vth$ is from the desired pitch $\vths$ computed in Step~1. If $\zth = 0$, the rocket points exactly where the position controller wants it to.
  \item $\zom$: how far the actual angular rate $\dot\vth$ is from the desired rate $\alpha_1$ computed in Step~2. If $\zom = 0$, the rocket is rotating at exactly the rate needed to drive $\zth$ to zero.
  \item $\zdel$: how far the actual nozzle angle $\delta$ is from the desired nozzle angle $\alpha_2$ computed in Step~3. If $\zdel = 0$, the nozzle is positioned exactly where needed to drive $\zom$ to zero.
\end{itemize}

Each error is only nonzero because the previous step's virtual control has not been fully achieved yet --- this chain structure is the essence of backstepping.

Each virtual control ($\vths$, $\alpha_1$, $\alpha_2$) is chosen to make the corresponding new term in $\dot V$ negative-definite, while passing a residual cross-coupling to the next step.

% =============================================================================
\section{Step 1 --- Virtual Control 1: Desired Pitch \texorpdfstring{$\vartheta^*$}{theta*}}
% =============================================================================

\textbf{Compute $\dot V_{\mathrm{pos}}$:}
\[
  \dot V_{\mathrm{pos}} = k_{px}e_x\dot x + \dot x\ddot x + k_{py}e_y\dot y + \dot y\ddot y
\]

Choose the desired accelerations:
\begin{align}
  A_x &= -k_{px}e_x - k_{dx}\dot x \notag \\
  A_y &= -k_{py}e_y - k_{dy}\dot y \tag{O1}
\end{align}

and require $(F/m)\sin\vths = A_x$, $(F/m)\cos\vths = A_y + g$. This gives:

\begin{remark}[Aerodynamics omitted from feedforward]
The desired accelerations (O1) account only for thrust, not for aerodynamic forces $X_b$, $Y_b$. Including them would make the computation circular: computing $\vths$ requires knowing $X_b, Y_b$, which depend on the angle of attack $\alpha = \vth - \mathrm{atan2}(\dot x, \dot y)$, which in turn depends on $\vths$. Instead, aerodynamic forces are treated as bounded disturbances --- they appear as extra terms when (T1)--(T2) are substituted into $\dot V_\mathrm{pos}$, and are absorbed into the ISS bound $D$ of Section~\ref{sec:iss}. This is justified because during the landing burn thrust dominates aerodynamics by roughly an order of magnitude, so the feedforward error is small and the feedback loop rejects the remainder.
\end{remark}

\begin{align}
  \sigma &= \clip\!\left(\frac{m\sqrt{A_x^2 + (A_y+g)^2}}{\Fmax},\, \sigma_{\min},\, 1\right) \tag{O2} \\[4pt]
  \vths &= \mathrm{atan2}(A_x,\, A_y + g) \tag{O3}
\end{align}

\textbf{Hover check:} $e_x = e_y = \dot x = \dot y = 0 \Rightarrow A_x = 0$, $A_y = 0$, $\sigma = mg/\Fmax$, $\vths = 0$. $\checkmark$

Now substitute (T1)--(T2) and (O1) into $\dot V_{\mathrm{pos}}$, applying A5: $\sin\vth \approx \sin\vths + \cos\vths\cdot \zth$, $\cos\vth \approx \cos\vths - \sin\vths\cdot \zth$. Using $(F/m)\sin\vths = A_x$ and $(F/m)\cos\vths = A_y + g$.

\textbf{Derivation of the $x$-channel term.} Substitute $\ddot x$ from (T1):
\[
  \dot x\ddot x + k_{px}e_x\dot x = \dot x\!\left(\frac{F\sin\vth + F\delta\cos\vth + X_b\sin\vth + Y_b\cos\vth}{m} + k_{px}e_x\right)
\]

Apply A5 to expand $F\sin\vth/m$ around $\vths$:
\[
  \frac{F\sin\vth}{m} \approx \frac{F\sin\vths}{m} + \frac{F\cos\vths}{m}\zth = A_x + \frac{F\cos\vths}{m}\zth
\]

Collecting terms:
\[
  \dot x\ddot x + k_{px}e_x\dot x = \dot x(A_x + k_{px}e_x) + \dot x\frac{F}{m}\cos\vths\cdot \zth + \delta\frac{F}{m}\cos\vth\cdot\dot x + \frac{\dot x(X_b\sin\vth + Y_b\cos\vth)}{m}
\]

With $A_x = -k_{px}e_x - k_{dx}\dot x$ from (O1), the first group gives $-k_{dx}\dot x^2$. The same expansion applied to the $y$-channel yields $-k_{dy}\dot y^2$ plus analogous residual terms.

The full aerodynamic contribution to $\dot V_{\mathrm{pos}}$ from both channels is:
\[
  \frac{\dot x(X_b\sin\vth + Y_b\cos\vth) + \dot y(X_b\cos\vth - Y_b\sin\vth)}{m}
  = \frac{X_b(\dot x\sin\vth + \dot y\cos\vth) + Y_b(\dot x\cos\vth - \dot y\sin\vth)}{m}
\]

The first group $\dot x\sin\vth + \dot y\cos\vth$ is the axial velocity $V_a$ (velocity component along the body axis), so $X_b V_a/m = -C_x q_\infty S_m V_a / m$, which is non-positive in forward flight ($V_a \geq 0$). The second group $\dot x\cos\vth - \dot y\sin\vth$ is the lateral velocity $V_n$, and $Y_b V_n/m = C_{y\alpha}\alpha q_\infty S_m V_n/m$. This is not sign-definite in general; it is absorbed into the ISS bound $D$ of Section~\ref{sec:iss} as an additional bounded term $\bar{d}_Y$.

Collecting the sign-definite terms:

\begin{equation}
  \dot V_{\mathrm{pos}} \leq -k_{dx}\dot x^2 - k_{dy}\dot y^2
  + \zth \underbrace{\frac{F}{m}(\dot x\cos\vths - \dot y\sin\vths)}_{\triangleq\, P}
  + \delta \underbrace{\frac{F}{m}(\dot x\cos\vth - \dot y\sin\vth)}_{\triangleq\, Q}
  + \bar{d}_Y
  \tag{DV1}
\end{equation}

where $\bar{d}_Y \geq 0$ is the bounded normal-force contribution. Note that $P$ uses $\vths$ (from the A5 linearisation) while $Q$ retains $\vth$ --- $Q$ is not linearised because it multiplies $\delta$, a small quantity by A4, and is carried as a residual to be cancelled in Step~4.

The term $P = (F/m)(\dot x\cos\vths - \dot y\sin\vths)$ is the \textbf{position-attitude coupling} arising from the tilt-velocity interaction. The term $Q = (F/m)(\dot x\cos\vth - \dot y\sin\vth)$ is the \textbf{nozzle-position coupling} arising from retaining $\delta$ in (T1)--(T2).

% =============================================================================
\section{Step 2 --- Virtual Control 2: Desired Angular Rate \texorpdfstring{$\alpha_1$}{alpha1}}
% =============================================================================

\textbf{Augment the Lyapunov function:} $V_2 = V_{\mathrm{pos}} + \tfrac{1}{2}\zth^2$.

\textbf{Compute $\dot\zth$:} $\dot\zth = \dot\vth - \dot\vths$.

From (DV1) we already have $\dot V_\mathrm{pos} \leq -k_{dx}\dot x^2 - k_{dy}\dot y^2 + \zth P + \delta Q + \bar d_Y$. Adding the new $\zth\dot\zth$ term and grouping the $\zth$ contributions, substituting $\dot\vth = \alpha_1 + \zom$:

\[
  \zth(P + \dot\zth) = \zth(P + \dot\vth - \dot\vths) = \zth(P + \alpha_1 + \zom - \dot\vths)
\]

\textbf{Design step.} Require $P + \alpha_1 - \dot\vths = -k_\vth \zth$:

\begin{equation}
  \boxed{\alpha_1 = \dot\vths - k_\vth \zth - \frac{F}{m}(\dot x\cos\vths - \dot y\sin\vths)}
  \tag{VC1}
\end{equation}

Then $\zth(P + \alpha_1 + \zom - \dot\vths) = -k_\vth \zth^2 + \zth\zom$.

\textbf{Result after Step 2:}
\begin{equation}
  \dot V_2 = -k_{dx}\dot x^2 - k_{dy}\dot y^2 - k_\vth \zth^2 + \zth\zom + \delta\cdot Q + \bar{d}_Y
  \tag{DV2}
\end{equation}

The correction $-P = -(F/m)(\dot x\cos\vths - \dot y\sin\vths)$ in $\alpha_1$ accounts for the fact that tilting the rocket changes the velocity components, which in turn affects the position error. Without this term, the $\zth\zom$ cross-coupling in $\dot V_2$ would not be cancelled and the Lyapunov inequality would have an uncompensated residual.

% =============================================================================
\section{Step 3 --- Virtual Control 3: Desired Nozzle Angle \texorpdfstring{$\alpha_2$}{alpha2}}
% =============================================================================

\textbf{Augment:} $V_3 = V_2 + \tfrac{1}{2}\zom^2$.

\textbf{Compute $\dot\zom$:} $\dot\zom = \ddot\vth - \dot\alpha_1 = g_2\,\delta + f_2 - \dot\alpha_1$.

\textbf{Time derivative:}
\begin{align*}
  \dot V_3 &= \dot V_2 + \zom\dot\zom \\
  &= -k_{dx}\dot x^2 - k_{dy}\dot y^2 - k_\vth \zth^2 + \zth\zom + \zom(g_2\,\delta + f_2 - \dot\alpha_1) + \delta\cdot Q + \bar{d}_Y
\end{align*}

Grouping all $\delta$-dependent terms, with $\Gamma \triangleq g_2\zom + Q$:
\begin{equation}
  \dot V_3 = -k_{dx}\dot x^2 - k_{dy}\dot y^2 - k_\vth \zth^2 + \zom(\zth + f_2 - \dot\alpha_1) + \delta\cdot\Gamma + \bar{d}_Y
  \tag{DV3}
\end{equation}

\textbf{Design step.} Write $\delta = \alpha_2 + \zdel$, so $\delta\cdot\Gamma = \alpha_2\cdot\Gamma + \zdel\cdot\Gamma$. Choose $\alpha_2$ so that the $\zom$ bracket gives $-k_\omega\zom^2$:
\[
  \zom(\zth + f_2 - \dot\alpha_1 + g_2\,\alpha_2) = -k_\omega\zom^2
\]

\begin{equation}
  \boxed{\alpha_2 = \frac{1}{g_2}\!\left(-k_\omega \zom - \zth - f_2 + \dot\alpha_1\right)}
  \tag{VC2}
\end{equation}

This requires $g_2 \neq 0$ --- guaranteed by Assumption A6.

\textbf{Result after Step 3:}
\begin{equation}
  \dot V_3 = -k_{dx}\dot x^2 - k_{dy}\dot y^2 - k_\vth \zth^2 - k_\omega\zom^2 + \zdel\cdot\Gamma + \alpha_2\cdot Q + \bar{d}_Y
  \tag{DV3'}
\end{equation}

Two residual terms are passed to Step 4:
\begin{itemize}
  \item $\zdel\cdot\Gamma = \zdel(g_2\zom + Q)$: coupling between nozzle-angle error and the rest of the system --- cancelled in Step~4.
  \item $\alpha_2\cdot Q$: coupling between the virtual nozzle reference and the nozzle-position cross-term. At equilibrium ($\dot x = \dot y = 0$), $Q = 0$, so this term vanishes. It is bounded during the transient.
\end{itemize}

% =============================================================================
\section{Step 4 --- Real Control: Nozzle Command \texorpdfstring{$\delta_{\mathrm{cmd}}$}{delta\_cmd}}
% =============================================================================

\textbf{Augment:}
\begin{equation}
  V = V_3 + \tfrac{1}{2}\zdel^2
  = \tfrac{1}{2}k_{px}e_x^2 + \tfrac{1}{2}\dot x^2 + \tfrac{1}{2}k_{py}e_y^2 + \tfrac{1}{2}\dot y^2 + \tfrac{1}{2}\zth^2 + \tfrac{1}{2}\zom^2 + \tfrac{1}{2}\zdel^2
  \tag{V4}
\end{equation}

\textbf{Compute $\dot\zdel$:}
\begin{equation}
  \dot\zdel = \dot\delta - \dot\alpha_2 = \frac{-\delta + \dcmd}{\tau_\delta} - \dot\alpha_2
  \tag{D4}
\end{equation}

\textbf{Time derivative:}
\begin{align*}
  \dot V &= \dot V_3 + \zdel\dot\zdel \\
  &= -k_{dx}\dot x^2 - k_{dy}\dot y^2 - k_\vth\zth^2 - k_\omega\zom^2 + \alpha_2\cdot Q
  + \zdel\!\left(\Gamma + \frac{-\delta + \dcmd}{\tau_\delta} - \dot\alpha_2\right) + \bar{d}_Y
\end{align*}

\textbf{Design step.} Choose $\dcmd$ to make the bracket multiplying $\zdel$ equal to $-k_\delta\zdel$:
\[
  \Gamma + \frac{-\delta + \dcmd}{\tau_\delta} - \dot\alpha_2 = -k_\delta\zdel
\]
Solving for $\dcmd$:

\begin{equation}
  \boxed{\dcmd = \delta + \tau_\delta\!\left(\dot\alpha_2 - g_2\zom - Q - k_\delta\zdel\right)}
  \tag{CL}
\end{equation}

\textbf{Verification of cancellation.} Substituting (CL) into the $\zdel$ bracket:
\begin{align*}
  \Gamma + \frac{-\delta + \dcmd}{\tau_\delta} - \dot\alpha_2
  &= (g_2\zom + Q) + (\dot\alpha_2 - g_2\zom - Q - k_\delta\zdel) - \dot\alpha_2
  = -k_\delta\zdel \quad \checkmark
\end{align*}

Apply saturation: $\dcmd = \clip(\dcmd,\,-\delta_{\max,\mathrm{cmd}},\,\delta_{\max,\mathrm{cmd}})$.

The term $-Q = -(F/m)(\dot x\cos\vth - \dot y\sin\vth)$ in (CL) feeds back the translational velocities into the nozzle command to compensate the position-nozzle coupling $\delta\cdot Q$ that appears in (DV1).

\textbf{$\dot V$ after substituting (CL):}
\begin{equation}
  \dot V = -k_{dx}\dot x^2 - k_{dy}\dot y^2 - k_\vth \zth^2 - k_\omega\zom^2 - k_\delta\zdel^2 + \alpha_2\cdot Q + \bar{d}_Y
  \tag{ISS1}
\end{equation}

% =============================================================================
\section{Stabilization Proof}
% =============================================================================

\subsection{Handling the \texorpdfstring{$\alpha_2 \cdot Q$}{alpha2 Q} Residual}

Before proceeding to the proof, we characterise the single term in (ISS1) that does not have a definite sign.

The coupling $Q = (F/m)(\dot x\cos\vth - \dot y\sin\vth)$ satisfies $|Q| \leq (F/m)\sqrt{\dot x^2 + \dot y^2}$. Applying Young's inequality with parameter $\epsilon > 0$:
\[
  |\alpha_2\cdot Q| \leq \frac{\epsilon}{2}\alpha_2^2 + \frac{1}{2\epsilon}Q^2 \leq \frac{\epsilon}{2}\alpha_2^2 + \frac{(F/m)^2}{2\epsilon}(\dot x^2 + \dot y^2)
\]

Choosing $\epsilon$ such that:
\begin{equation}
  \epsilon \geq \frac{(F/m)^2}{2\min(k_{dx}, k_{dy})}
  \tag{C-eps}
\end{equation}

the $Q^2$ residual is absorbed by the velocity damping terms, and (ISS1) becomes:

\begin{equation}
  \dot V \leq -\!\left(k_{dx} - \frac{(F/m)^2}{2\epsilon}\right)\!\dot x^2 - \!\left(k_{dy} - \frac{(F/m)^2}{2\epsilon}\right)\!\dot y^2
  - k_\vth \zth^2 - k_\omega\zom^2 - k_\delta\zdel^2 + \frac{\epsilon}{2}\alpha_2^2
  \tag{ISS2}
\end{equation}

Define $k_{dx}' = k_{dx} - (F/m)^2/(2\epsilon) > 0$ and $k_{dy}' = k_{dy} - (F/m)^2/(2\epsilon) > 0$ (both positive by (C-eps)). The remaining term $\tfrac{\epsilon}{2}\alpha_2^2$ is positive but vanishes at equilibrium since $\alpha_2 \to 0$ when $\zom, \zth, f_2 \to 0$ in (VC2). Its treatment is in Section~\ref{sec:uub}.

\subsection{Boundedness of All Signals}
\label{sec:uub}

\textbf{Note on the $\tfrac{\epsilon}{2}\alpha_2^2$ term.} From (VC2), $\alpha_2 = \tfrac{1}{g_2}(-k_\omega\zom - \zth - f_2 + \dot\alpha_1)$, so $\alpha_2$ is a function of $\zom$, $\zth$, $f_2$, and $\dot\alpha_1$, all of which are components of the state vector covered by $V$. Taking the absolute value and using $g_2 \geq g_{2,\min} > 0$ (Assumption A6):
\[
  |\alpha_2| \leq \frac{1}{g_{2,\min}}\!\left(k_\omega|\zom| + |\zth| + |f_2| + |\dot\alpha_1|\right)
\]

Since $V \geq \tfrac{1}{2}\zth^2$ and $V \geq \tfrac{1}{2}\zom^2$:
\begin{align}
  |\zth| &\leq \sqrt{2V} \\
  |\zom| &\leq \sqrt{2V}
\end{align}

\textbf{Boundedness of $|f_2|$.} The aerodynamic moment $f_2 = C_{m\alpha}\,\alpha\,q_\infty S_m l/J$ depends on airspeed $V = \sqrt{\dot x^2 + \dot y^2}$ (bounded by physical limits of the scenario) and angle of attack $\alpha = \vth - \mathrm{atan2}(\dot x, \dot y)$. The pitch $\vth$ is bounded via $\vth = \vths + \zth$, where $\vths = \mathrm{atan2}(A_x, A_y+g)$ depends only on the bounded position errors $e_x, \dot x, e_y, \dot y$ (covered by $V_\mathrm{pos} \leq V$). Therefore:
\[
  |f_2| \leq \frac{C_{m\alpha}\,|\alpha|_{\max}\,q_{\infty,\max}\,S_m\,l}{J} \triangleq \bar f_2 < \infty
\]

\textbf{Boundedness of $|\dot\alpha_1|$.} Differentiating (VC1): $\dot\alpha_1 = \ddot\vths - k_\vth\dot\zth - \dot P$. Each term depends on positions, velocities, and their derivatives --- all bounded by $V$ and the physical equations of motion. Therefore $|\dot\alpha_1| \leq \bar\alpha_1 < \infty$.

Combining all four bounds, there exists a constant $C_{\alpha_2}$ depending only on $V(0)$ and the gains such that:
\[
  \frac{\epsilon}{2}\alpha_2^2 \leq \frac{\epsilon}{2}C_{\alpha_2}^2 \triangleq \mu
\]

Substituting into (ISS2):
\begin{equation}
  \dot V \leq -k_{dx}'\,\dot x^2 - k_{dy}'\,\dot y^2 - k_\vth \zth^2 - k_\omega\zom^2 - k_\delta\zdel^2 + \mu
  \tag{ISS3}
\end{equation}

\begin{theorem}[UUB]
Under the control law (CL), virtual controls (VC1)--(VC2), and Assumptions A1--A7, if gains satisfy (C-eps) strictly, then all signals $(e_x, \dot x, e_y, \dot y, \zth, \zom, \zdel)$ are \textbf{uniformly ultimately bounded}: they enter and remain in the compact set $\Omega = \{V \leq \mu/c\}$ in finite time, where:
\[
  c = \min(k_{dx}',\, k_{dy}',\, k_\vth,\, k_\omega,\, k_\delta)
\]
The transient bound is $V(t) \leq \max(V(0),\, \mu/c) \triangleq V_{\max}$. Individual bounds:
\begin{align}
  |e_x(t)| &\leq \sqrt{\frac{2V_{\max}}{k_{px}}} \triangleq B_{e_x} \\
  |\dot x(t)| &\leq \sqrt{2V_{\max}} \triangleq B_{\dot x}
\end{align}
and analogously for $e_y$, $\dot y$, $\zth$, $\zom$, $\zdel$. All bounds depend only on $V(0)$, $\mu$, and the gains.
\end{theorem}

\textbf{Consequence.} Since $(\zth, \zom, \zdel)$ are bounded and the virtual controls (VC1)--(VC2) are continuous functions of bounded signals, $\vth$, $\dot\vth$, $\delta$ and $\alpha_2$ are all bounded. In particular, the singularity $g_2 = 0$ is never reached under A6.

\begin{remark}[Two-phase stability picture]
The overall stability result has two phases:
\begin{itemize}
  \item \textbf{Phase 1 --- transient ($\alpha_2 \neq 0$, $\mu > 0$):} While the system is away from equilibrium, $\alpha_2 \neq 0$ and the residual $\mu = \tfrac{\epsilon}{2}C_{\alpha_2}^2 > 0$ from the $\alpha_2 \cdot Q$ cross-term is nonzero. Additionally, clipping of $\alpha_1$, $\alpha_2$, $\dcmd$ may be active (e.g.\ during the first 1--2 s with a $30^{\circ}$ initial tilt), contributing the extra disturbance $D'$ of Section~\ref{sec:iss}. In this phase only UUB holds: all errors are bounded and remain in the ball $\Omega = \{V \leq D'/c\}$. The system does not diverge, but exact convergence to zero is not guaranteed.
  \item \textbf{Phase 2 --- steady descent ($\alpha_2 \to 0$, $\mu \to 0$):} As the system approaches equilibrium, $\alpha_2 \to 0$ so $\mu \to 0$, and clipping becomes inactive so $D' \to 0$. The stronger results of Sections~\ref{sec:barbalat} and~\ref{sec:lasalle} apply: Barbalat gives $\dot x, \dot y, \zth, \zom, \zdel \to 0$, and LaSalle gives $e_x, e_y \to 0$.
\end{itemize}
In short: the controller is practically stable throughout, and asymptotically convergent once $\alpha_2$ is small and clipping is inactive.
\end{remark}

\subsection{Convergence of Velocities and Attitude Errors --- Barbalat's Lemma}
\label{sec:barbalat}

\begin{remark}[Scope]
The Barbalat argument below is an idealised analysis that assumes $\mu = 0$ from the start --- i.e., the $\alpha_2 \cdot Q$ cross-term is neglected. This is a standard simplification used to establish the convergence structure of the controller. In reality $\mu = \tfrac{\epsilon}{2}C_{\alpha_2}^2 > 0$ whenever $\alpha_2 \neq 0$, so the true result is UUB (Section~\ref{sec:uub}), not exact convergence to zero. The Barbalat argument is nonetheless useful because it shows what the system would do in the absence of the residual, and because $\mu$ is small when $\alpha_2$ is small (i.e., when the nozzle is near its equilibrium position).
\end{remark}

\begin{theorem}[Idealised, $\mu = 0$]
If the $\alpha_2 \cdot Q$ residual is neglected, then $\dot x(t) \to 0$, $\dot y(t) \to 0$, $\zth(t) \to 0$, $\zom(t) \to 0$, $\zdel(t) \to 0$ as $t \to \infty$.
\end{theorem}

\begin{proof}
We show the argument for $\dot x$; the others follow identically.

\textit{Step 1 --- Integrability.} With $\mu = 0$, (ISS3) gives $\dot V \leq -k_{dx}'\,\dot x^2$. Integrating from $0$ to $T$:
\[
  \int_0^T k_{dx}'\,\dot x^2(\tau)\,d\tau \leq V(0) - V(T) \leq V(0) < \infty
\]
Taking $T \to \infty$: $\dot x^2 \in L^2([0,\infty))$.

\textit{Step 2 --- Uniform continuity.} From Section~\ref{sec:uub}, $|\dot x| \leq \sqrt{2V_{\max}}$. From (T1), $\ddot x$ is a continuous function of bounded signals, so $\tfrac{d}{dt}(\dot x^2) = 2\dot x\,\ddot x$ is bounded, meaning $\dot x^2$ is uniformly continuous.

\textit{Step 3 --- Barbalat's Lemma.} A non-negative uniformly continuous function with finite $L^2$ integral converges to zero:
\[
  \dot x^2(t) \to 0 \implies \dot x(t) \to 0 \quad \text{as } t \to \infty. \qquad \square
\]
\end{proof}

The same argument applies to $\dot y$, $\zth$, $\zom$, and $\zdel$ using the corresponding negative-definite terms in (ISS3).

\subsection{Convergence of Position Errors --- LaSalle's Invariance Principle}
\label{sec:lasalle}

\begin{theorem}[Idealised, same scope as Section~\ref{sec:barbalat}]
If $\mu = 0$, then $e_x(t) \to 0$ and $e_y(t) \to 0$ as $t \to \infty$. Under the true $\mu > 0$, position errors are bounded by $|e_x| \leq B_{e_x}$ from Section~\ref{sec:uub}; exact convergence to zero is not claimed.
\end{theorem}

\begin{proof}
\textbf{Applicability of LaSalle.} The closed-loop system is \textbf{autonomous} --- the right-hand side $f(s)$ depends only on the state $s$, not explicitly on $t$ (the controller is a pure state-feedback law with no external time-varying signals). LaSalle's invariance principle (Khalil, Theorem 4.4) therefore applies directly. The set $\{\dot V = 0\}$ contains states where $\dot x = \dot y = \zth = \zom = \zdel = 0$ but $e_x, e_y$ may be non-zero, so we must identify the largest invariant subset of $\{\dot V = 0\}$. All trajectories remain in the compact sublevel set $\{V \leq V_{\max}\}$ (Section~\ref{sec:uub}), satisfying the compactness requirement of the theorem.

Define the limiting set:
\[
  \mathcal{S} = \bigl\{(e_x,\, \dot x,\, e_y,\, \dot y,\, \zth,\, \zom,\, \zdel) : \dot x = \dot y = \zth = \zom = \zdel = 0\bigr\}
\]

\textbf{In $\mathcal{S}$} the following hold simultaneously:
\begin{itemize}
  \item $\dot x = \dot y = 0 \Rightarrow \ddot x = \ddot y = 0$; also $Q = 0$.
  \item $\zth = 0 \Rightarrow \vth = \vths = \mathrm{atan2}(A_x,\, A_y+g)$.
  \item $\zom = 0 \Rightarrow \dot\vth = \alpha_1$; with $\dot x = \dot y = 0$ and $\zth = 0$: $\alpha_1 = \dot\vths - 0 - 0 = \dot\vths$, so $\dot\vth = \dot\vths$.
  \item $\zdel = 0 \Rightarrow \delta = \alpha_2$.
\end{itemize}

Substituting $\dot x = \dot y = 0$ into (T1) (aerodynamic terms vanish since $V = 0$):
\[
  m\ddot x = F\sin\vths
\]
By construction (O3), $(F/m)\sin\vths = A_x$, so $F\sin\vths = mA_x$, giving $m\ddot x = mA_x$.

But $\ddot x = 0$ in $\mathcal{S}$, so:
\[
  0 = A_x = -k_{px}e_x - k_{dx}\underbrace{\dot x}_{=0} = -k_{px}e_x \implies e_x = 0
\]
The same argument applied to (T2) gives $e_y = 0$.

Therefore the largest invariant subset of $\mathcal{S}$ is $\{e_x = e_y = 0\}$, and $(e_x(t),\, e_y(t)) \to 0$. $\square$
\end{proof}

% =============================================================================
\subsection{Stability Under Saturation --- ISS Analysis}
\label{sec:iss}
% =============================================================================

Section~\ref{sec:uub} (UUB) holds regardless of clipping. The idealised convergence proofs in Sections~\ref{sec:barbalat}--\ref{sec:lasalle} additionally assume that $\alpha_1$, $\alpha_2$, and $\dcmd$ are delivered exactly. In the real implementation all three are clipped:

\begin{align}
  \alpha_{1,\mathrm{sat}} &= \clip(\alpha_{1,\mathrm{raw}},\,-\alpha_{1,\max},\,\alpha_{1,\max}) \\
  \alpha_{2,\mathrm{sat}} &= \clip(\alpha_{2,\mathrm{raw}},\,-\dmax,\,\dmax) \\
  \delta_{\mathrm{cmd},\mathrm{sat}} &= \clip(\dcmd,\,-\delta_{\max,\mathrm{cmd}},\,\delta_{\max,\mathrm{cmd}})
\end{align}

This section proves that the closed-loop system remains \textbf{input-to-state stable (ISS)} with respect to the saturation errors, giving a rigorous global result without assuming saturation never occurs.

\subsubsection{Saturation Errors as Bounded Disturbances}

Each saturation is modelled as an additive bounded disturbance. All three disturbances are uniformly bounded regardless of state.

\textbf{$\alpha_1$ saturation:}
\begin{equation}
  d_\vartheta = \alpha_{1,\mathrm{sat}} - \alpha_{1,\mathrm{raw}}, \quad |d_\vartheta| \leq 2\alpha_{1,\max} \triangleq \bar{d}_\vartheta
\end{equation}
This enters the $\zom$ dynamics as $-d_\vartheta$ (since $\zom = \dot\vth - \alpha_1$ and $\alpha_1$ is clipped).

\textbf{$\alpha_2$ saturation:}
\begin{equation}
  d_\omega = \alpha_{2,\mathrm{sat}} - \alpha_{2,\mathrm{raw}}, \quad |d_\omega| \leq 2\dmax
\end{equation}
This propagates through the $\zom$ dynamics with gain $g_2$, giving bounded disturbance $|g_2 d_\omega| \leq g_{2,\max}\cdot 2\dmax \triangleq \bar{d}_\omega$.

\textbf{$\dcmd$ saturation:}
\begin{equation}
  d_\delta = \delta_{\mathrm{cmd},\mathrm{sat}} - \dcmd, \quad |d_\delta| \leq 2\delta_{\max,\mathrm{cmd}}
\end{equation}
This enters the $\zdel$ dynamics divided by $\tau_\delta$, giving bounded disturbance $|d_\delta/\tau_\delta| \leq 2\delta_{\max,\mathrm{cmd}}/\tau_\delta \triangleq \bar{d}_\delta$.

\subsubsection{Perturbed Error Dynamics}

Substituting $\alpha_{2,\mathrm{sat}} = \alpha_{2,\mathrm{raw}} + d_\omega$ and using the design choice (VC2):
\[
  \dot\zom = -k_\omega\zom - \zth + g_2\zdel + g_2\,d_\omega - d_\vartheta
\]

Substituting $\delta_{\mathrm{cmd},\mathrm{sat}} = \dcmd + d_\delta$ and using (CL):
\[
  \dot\zdel = -k_\delta\zdel - g_2\zom - Q + \frac{d_\delta}{\tau_\delta}
\]

\subsubsection{Lyapunov Derivative Under Disturbances}

The $z_\omega$ and $z_\delta$ contributions to $\dot V$ are:
\[
  \zom\,\dot\zom = -k_\omega\zom^2 - \zth\zom + g_2\zom\zdel + g_2\,d_\omega\,\zom - d_\vartheta\zom
\]
\[
  \zdel\,\dot\zdel = -k_\delta\zdel^2 - g_2\zom\zdel - Q\,\zdel + \frac{d_\delta}{\tau_\delta}\,\zdel
\]

The $g_2\zom\zdel$ terms cancel exactly. The cross-terms are bounded via Young's inequality:
\begin{itemize}
  \item $|g_2 d_\omega\zom| \leq \dfrac{k_\omega}{8}\zom^2 + \dfrac{2g_2^2\bar{d}_\omega^2}{k_\omega}$ \quad (with $\lambda_1 = k_\omega/4$)
  \item $\left|\dfrac{d_\delta}{\tau_\delta}\zdel\right| \leq \dfrac{k_\delta}{16}\zdel^2 + \dfrac{4\bar{d}_\delta^2}{k_\delta\tau_\delta^2}$ \quad (with $\lambda_2 = k_\delta/8$)
  \item $|d_\vartheta\zom| \leq \dfrac{k_\omega}{8}\zom^2 + \dfrac{2\bar{d}_\vartheta^2}{k_\omega}$ \quad (with $\lambda_5 = k_\omega/4$)
\end{itemize}

These reduce the $\zom^2$ coefficient by $k_\omega/8 + k_\omega/8 = k_\omega/4$ total and the $\zdel^2$ coefficient by $k_\delta/16$, keeping both positive. The $\zth\zom$ cross-term (from (DV2)) is bounded with $\lambda_4 = k_\vth/2$:
\[
  |\zth\zom| \leq \frac{k_\vth}{4}\zth^2 + \frac{1}{k_\vth}\zom^2
\]
provided $k_\omega > 1/k_\vth + 1$. The $Q\zdel$ term is absorbed into the velocity terms via (C-eps) extended to cover $\lambda_3 = k_\delta/4$.

Collecting all terms:
\begin{equation}
  \dot V \leq -k_{dx}''\,\dot x^2 - k_{dy}''\,\dot y^2 - \frac{k_\vth}{2}\zth^2 - \frac{k_\omega}{4}\zom^2 - \frac{k_\delta}{8}\zdel^2 + D'
  \tag{ISS4}
\end{equation}

where the total disturbance bound is:
\begin{equation}
  D' = \mu + \frac{2g_2^2\bar{d}_\omega^2}{k_\omega} + \frac{2\bar{d}_\vartheta^2}{k_\omega} + \frac{4\bar{d}_\delta^2}{k_\delta\tau_\delta^2}
\end{equation}

\subsubsection{UUB Conclusion}

\begin{theorem}[UUB under saturation]
Under Assumptions A1--A7, gain condition (C-eps), and $k_\omega > 1/k_\vth + 1$, the closed-loop system with saturated controls is \textbf{uniformly ultimately bounded}: all states enter and remain in $\mathcal{B} = \{V \leq D/c\}$, where:
\begin{align}
  c &= \min\!\left(k_{dx}'',\, k_{dy}'',\, \frac{k_\vth}{2},\, \frac{k_\omega}{4},\, \frac{k_\delta}{8}\right) \\
  D &= D'
\end{align}
The residual bounds on attitude errors are:
\begin{align}
  |z_\omega|_{\infty} &\leq \sqrt{\frac{2D}{c\,k_\omega}} \\
  |z_\delta|_{\infty} &\leq \sqrt{\frac{2D}{c\,k_\delta}}
\end{align}
\end{theorem}

\subsubsection{Quantitative Bound for the Project Parameters}

With parameters from \texttt{configs/backstepping.yaml}:
\begin{itemize}
  \item $g_{2,\max} = \Fmax\,l_{cp}/J = 191763 \times 10.5 / 318196.65 \approx 6.33$ rad/s$^2$/rad
  \item $\alpha_{1,\max} = 1.5$ rad/s, so $\bar{d}_\vartheta = 3.0$ rad/s
  \item $\dmax = 15^{\circ} = 0.262$ rad, so $\bar{d}_\omega = 2 \times 6.33 \times 0.262 \approx 3.31$ rad/s$^2$
  \item $\delta_{\max,\mathrm{cmd}} = 0.262$ rad, $\tau_\delta = 0.05$ s, so $\bar{d}_\delta \approx 10.47$ rad/s
\end{itemize}

With $k_\omega = 5$, $k_\delta = 15$, $\tau_\delta = 0.05$ s, the three contributions to $D'$ are:
\begin{align}
  \frac{2g_{2,\max}^2\bar{d}_\omega^2}{k_\omega} &\approx 175.8 \\
  \frac{2\bar{d}_\vartheta^2}{k_\omega} &\approx 3.6 \\
  \frac{4\bar{d}_\delta^2}{k_\delta\tau_\delta^2} &\approx 11697
\end{align}

The $\dcmd$ clipping term dominates --- it is amplified by $1/\tau_\delta^2 = 400$. Total $D' \approx \mu + 11876$.

With $c = \min(k_\vth/2,\, k_\omega/4,\, k_\delta/8) = \min(3.0,\, 1.25,\, 1.875) = 1.25$:
\[
  |\zom|_\infty \leq \sqrt{\frac{2 \times 11876}{1.25 \times 5}} \approx 61.6\ \text{rad/s} \approx 3532^{\circ}/\text{s}
\]

This absurdly large bound is a direct consequence of the $1/\tau_\delta^2$ amplification --- the ISS framework assumes $\dcmd$ clipping is permanently active, which never happens in practice. In the simulation, $\dcmd$ clipping clears within 1--2 s; after that $d_\delta = 0$ exactly, $D' \to \mu$, and the system converges asymptotically. The bound is technically correct but extremely conservative.

This large bound assumes all three clipping saturations are \textbf{permanently active simultaneously}. In the simulation, $\alpha_1$ clipping clears within 0.1 s and $\alpha_2$ clipping within 1--2 s; after that $d_\vartheta = d_\omega = d_\delta = 0$ exactly. As the system then converges, $\alpha_2 \to 0$ so $\mu \to 0$, and the results of Sections~\ref{sec:barbalat}--\ref{sec:lasalle} apply, giving true asymptotic convergence to zero.

% =============================================================================
\section{Full Algorithm}
% =============================================================================

\textbf{State:} $s = [x, y, \dot x, \dot y, \vth, \dot\vth, \delta]^T$. Filter state: $\alpha_2^f$ (command filter --- introduced in Block~E to avoid differentiating $\alpha_2$ analytically). Derivative filter states: $\dot\vths_{\mathrm{filt}}$, $\dot\alpha_{1,\mathrm{filt}}$ (ODE states, see Section~\ref{sec:impl}).

\textbf{Block A --- Step 1: Virtual control 1 ($\vths$, $\sigma$)}

\begin{lstlisting}
1.  e_x = x - x_d,   e_y = y
2.  A_x = -k_px * e_x - k_dx * xdot
3.  A_y = -k_py * e_y - k_dy * ydot
4.  a_vert = max(A_y + g, |A_x| / tan(phi_lim), a_min)   [safety clamp]
5.  A_y    = a_vert - g
6.  A_x    = clip(A_x, -a_vert*tan(phi_lim), a_vert*tan(phi_lim))
7.  T_des  = sqrt(A_x^2 + a_vert^2)
8.  sigma  = clip(m * T_des / F_max, sigma_min, 1)
9.  F      = sigma * F_max
10. theta* = atan2(A_x, a_vert)
11. theta*_dot ~= filtered analytically (see Section 13)
\end{lstlisting}

\begin{remark}[Safety clamping, steps 4--6]
The theoretical design requires only $a_{\min} = 0$ for the Lyapunov proof; the implementation enforces $a_{\min} = 0.5$ m/s$^2$ and the tilt limit $\phi_{\lim} = 30^{\circ}$ to keep the rocket from commanding full free-fall or excessive tilt. The clamp on $A_x$ ensures $|\vths| \leq \phi_{\lim}$ exactly. These modifications introduce a bounded error in (O1) handled by Section~\ref{sec:iss}.
\end{remark}

\textbf{Block B --- Aerodynamics}

\begin{lstlisting}
12. V_spd = sqrt(xdot^2 + ydot^2)
13. alpha_aoa = (V_spd < V_min) ? theta : theta - atan2(xdot, ydot)
14. q_inf = 0.5 * rho * V_spd^2
15. f2 = C_ma * alpha_aoa * q_inf * Sm * l / J
16. g2 = F * l_cp / J
\end{lstlisting}

\textbf{Block C --- Step 2: Virtual control 2 ($\alpha_1$)}

\begin{lstlisting}
17. z_theta = theta - theta*
18. alpha1  = theta*_dot - k_phi * z_theta
19. z_omega = theta_dot - alpha1
20. alpha1_dot ~= (alpha1(t) - alpha1(t-dt)) / dt
\end{lstlisting}

\begin{remark}[P-coupling omission]
The theoretical formula (VC1) includes $-(F/m)(\dot x\cos\vths - \dot y\sin\vths)$. In practice this reaches 60+ rad/s at high entry velocities and permanently saturates $\alpha_1$. The simplified form $\alpha_1 = \dot\vths - k_\vth\zth$ is used instead. The omitted P term is a bounded disturbance covered by Section~\ref{sec:iss}.
\end{remark}

\begin{remark}[Filtered derivatives]
The implementation propagates $\dot\vths$ and $\dot\alpha_1$ as ODE states driven by first-order filters toward their instantaneous analytical values (bandwidth $N = 50$ rad/s). The lag error is bounded by $O(1/N)$ and constitutes a bounded disturbance within the ISS framework.
\end{remark}

\begin{remark}[$\alpha_1$ saturation]
The implementation clips $\alpha_1$ to $\pm\alpha_{1,\max}$ (default 1.5 rad/s). This introduces a bounded disturbance in the $\zom$ dynamics analogous to the $\alpha_2$ saturation analysed in Section~\ref{sec:iss}, with $\bar{d}_{\alpha_1} = 2\alpha_{1,\max}$.
\end{remark}

\textbf{Block D --- Step 3: Virtual control 3 ($\alpha_2$)}

\begin{lstlisting}
21. alpha2_raw = (1/g2) * (-k_omega * z_omega - z_theta - f2 + alpha1_dot)
22. alpha2_sat = clip(alpha2_raw, -delta_max, delta_max)
\end{lstlisting}

\textbf{Block E --- Command filter}

\begin{lstlisting}
23. alpha2f_dot = (alpha2_sat - alpha2f) / tau_f
24. alpha2f    += alpha2f_dot * dt    [or integrate inside ODE]
25. z_delta     = delta - alpha2f
\end{lstlisting}

\begin{remark}[Command filtering]
The Lyapunov derivation defines $\zdel = \delta - \alpha_2$ and uses $\dot\alpha_2$ in (CL). In the implementation, $\alpha_2$ is replaced by its filtered version $\alpha_2^f$, so $\zdel = \delta - \alpha_2^f$ and $\dot\alpha_2^f$ enters (CL). The filter error $e_f = \alpha_2^f - \alpha_2$ satisfies $\dot e_f = -e_f/\tau_f + (\alpha_{2,\mathrm{sat}} - \alpha_2)/\tau_f$, which is bounded. Command filtering (Farrell et al., 2005) avoids differentiating $\alpha_2$ analytically; the resulting filter error introduces an additional bounded disturbance of order $O(\tau_f)$ absorbed by Section~\ref{sec:iss}.
\end{remark}

\textbf{Block F --- Step 4: Real control ($\dcmd$)}

\begin{lstlisting}
26. delta_cmd = delta + tau_delta * (alpha2f_dot - g2*z_omega - k_delta*z_delta)
27. delta_cmd = clip(delta_cmd, -delta_max_cmd, delta_max_cmd)
\end{lstlisting}

\begin{remark}[Q-coupling omission]
The theoretical formula (CL) includes $-Q = -(F/m)(\dot x\cos\vth - \dot y\sin\vth)$. This term reaches 50+ rad/s at high entry velocities and, multiplied by $\tau_\delta$, produces nozzle commands far exceeding $\delta_{\max,\mathrm{cmd}}$. The term is omitted in the implementation; the bounded error is covered by Section~\ref{sec:iss}.
\end{remark}

\textbf{Output:} $[\sigma,\ \dcmd]$

% =============================================================================
\section{Gain Conditions and Tuning Guidelines}
% =============================================================================

\subsection*{Necessary conditions}
\[
  k_{px}, k_{dx}, k_{py}, k_{dy}, k_\vth, k_\omega, k_\delta > 0, \quad \sigma_{\min} > 0, \quad \tau_f > 0
\]

\subsection*{Critical damping of position loop}
\begin{align}
  k_{dx} &\geq 2\sqrt{k_{px}} \\
  k_{dy} &\geq 2\sqrt{k_{py}}
\end{align}

\subsection*{Attitude loop bandwidth}

No formal time-scale separation is required, but faster attitude convergence reduces the time during which the $\alpha_2\cdot Q$ residual is active. A practical guideline:
\begin{align}
  \sqrt{k_\vth\, g_{2,\mathrm{hov}}} &\geq 3\,\sqrt{k_{px}} \\
  g_{2,\mathrm{hov}} &= \frac{\sigma_{\mathrm{hov}} \Fmax l_{cp}}{J}
\end{align}

\subsection*{Cross-term damping condition}

From the $\zth\zom$ bounding step in Section~\ref{sec:iss}:
\[
  k_\omega > \frac{1}{k_\vth} + 1
\]

For the project gains $k_\vth = 6$: $k_\omega > 1/6 + 1 \approx 1.17$. Current value $k_\omega = 5$ satisfies this with large margin.

\subsection*{Residual $\alpha_2 \cdot Q$ condition (C-eps)}

From Section~10.1, the Young's inequality bound on $|\alpha_2 \cdot Q|$ introduces the term $(F/m)^2(\dot x^2 + \dot y^2)/(2\epsilon)$, which must be absorbed by the velocity damping. This requires:
\[
  \frac{(F/m)^2}{2\epsilon} \leq \min(k_{dx},\, k_{dy}) \quad\Longleftrightarrow\quad \epsilon \geq \frac{(F/m)^2}{2\min(k_{dx},\, k_{dy})}
\]

For the project parameters, $(F/m)_{\max} = \Fmax/m = 191763/13000 \approx 14.75$ m/s$^2$, so $\epsilon \geq 14.75^2/(2\times 5.0) \approx 21.7$. Since $\epsilon$ is a free parameter in Young's inequality (not a gain to tune), this condition is always satisfiable.

\subsection*{Command filter time constant}
\[
  \frac{\tau_\delta}{10} \lesssim \tau_f \lesssim \frac{\tau_\delta}{5}
\]

% =============================================================================
\section{Implementation Notes}
\label{sec:impl}
% =============================================================================

\subsection*{Ten-state ODE}

\begin{lstlisting}[language=Python]
state = [x, y, vx, vy, phi, omega, delta, alpha2f, ts_dot, a1_dot]
\end{lstlisting}

The last two states are first-order filtered derivatives: \texttt{ts\_dot} tracks $\dot\vths$ and \texttt{a1\_dot} tracks $\dot\alpha_1$, driven analytically inside the ODE.

\subsection*{Filtered derivatives}

$\dot\vths$ and $\dot\alpha_1$ are needed by (VC2) and (CL). The implementation propagates them as ODE states driven by first-order filters toward their analytical instantaneous values:

\begin{lstlisting}[language=Python]
ts_dot_instant = ((A_y+g)*dAx_dt - A_x*dAy_dt) / max(A_x**2+(A_y+g)**2, a_vert**2)
ts_dot_instant = clip(ts_dot_instant, -alpha1_max, alpha1_max)
a1_dot_instant = ts_dot_instant - k_phi * (omega - ts_dot)

d/dt(ts_dot) = N_deriv * (ts_dot_instant - ts_dot)   # N_deriv = 50 rad/s
d/dt(a1_dot) = N_deriv * (a1_dot_instant - a1_dot)
\end{lstlisting}

This avoids differentiating $\vths$ numerically (which is noisy), at the cost of a lag error bounded by $O(1/N_{\mathrm{deriv}})$.

\subsection*{Angle wrapping}

\begin{lstlisting}[language=Python]
z_phi = (phi - theta_star + np.pi) % (2 * np.pi) - np.pi
\end{lstlisting}

\subsection*{Zero-airspeed guard}

\begin{lstlisting}[language=Python]
if V_spd < V_min:   # V_min ~= 0.1 m/s
    alpha_aoa, q_inf = phi, 0.0
else:
    alpha_aoa = phi - np.arctan2(vx, vy)
    q_inf = 0.5 * rho * V_spd**2
\end{lstlisting}

\subsection*{Minimum throttle guard}

\begin{lstlisting}[language=Python]
sigma = max(sigma, sigma_min)
F = sigma * F_max
g2 = F * l_cp / J   # guaranteed > 0
\end{lstlisting}

\subsection*{Python variable mapping}

\begin{center}
\begin{tabular}{ll}
\toprule
Symbol & Python variable \\
\midrule
$\vth$ & \texttt{phi} \\
$\dot\vth$ & \texttt{omega} \\
$\vths$ & \texttt{theta\_star} \\
$\delta$ & \texttt{delta} \\
$\dcmd$ & \texttt{delta\_cmd} \\
$\alpha_2^f$ & \texttt{alpha2f} \\
$\zth,\, \zom,\, \zdel$ & \texttt{z\_phi, z\_omega, z\_delta} \\
$\alpha_1,\, \alpha_2$ & \texttt{alpha1, alpha2} \\
$P,\, Q$ & \texttt{P\_coupling, Q\_coupling} \\
$f_2,\, g_2$ & \texttt{f2, g2} \\
$\sigma$ & \texttt{sigma} \\
$k_\vth,\, k_\omega,\, k_\delta$ & \texttt{k\_phi, k\_omega, k\_delta} \\
$k_{px}, k_{dx}, k_{py}, k_{dy}$ & \texttt{k\_px, k\_dx, k\_py, k\_dy} \\
$\tau_\delta,\, \tau_f$ & \texttt{tau\_delta, tau\_filter} \\
\bottomrule
\end{tabular}
\end{center}

% =============================================================================
\section*{References}
% =============================================================================

\begin{enumerate}
  \item Khalil, H.\,K. (2002). \textit{Nonlinear Systems} (3rd ed.). Prentice Hall. --- Lyapunov stability theory, LaSalle's invariance principle, Young's inequality applications.
  \item Krsti\'{c}, M., Kanellakopoulos, I., \& Kokotovi\'{c}, P. (1995). \textit{Nonlinear and Adaptive Control Design}. Wiley. --- Backstepping for strict-feedback systems (Chapter 2); composite Lyapunov functions; residual cross-term handling.
  \item Farrell, J.\,A., Sharma, M., \& Polycarpou, M. (2005). Backstepping-based flight control with adaptive function approximation. \textit{AIAA Journal of Guidance, Control, and Dynamics}, 28(6), 1089--1102. --- Command filtering; UUB under filter error.
  \item Slotine, J.-J.\,E., \& Li, W. (1991). \textit{Applied Nonlinear Control}. Prentice Hall. --- Robustness of Lyapunov designs; cross-term bounding techniques.
  \item Wie, B. (1998). \textit{Space Vehicle Dynamics and Control}. AIAA. --- TVC torque derivation; nozzle actuator modelling.
  \item Sontag, E.\,D. (1989). Smooth stabilization implies coprime factorization. \textit{IEEE Trans.\ Autom.\ Control}, 34(4), 435--443. --- Original ISS definition and the comparison-lemma decay bound used in Section~\ref{sec:iss}.
  \item Sontag, E.\,D., \& Wang, Y. (1995). On characterizations of the input-to-state stability property. \textit{Systems \& Control Letters}, 24(5), 351--359. --- ISS Lyapunov characterisation; the $\dot V \leq -cV + D$ criterion.
\end{enumerate}

\end{document}
